\documentclass[12pt,a4paper]{article}

% ------------------------------------------------
% Pacotes
% ------------------------------------------------
\usepackage[utf8]{inputenc}
\usepackage[T1]{fontenc}
\usepackage[brazil]{babel}
\usepackage{graphicx}
\usepackage[table]{xcolor}
\usepackage{array}
\usepackage{tabularx}
\usepackage{geometry}
\usepackage{fancyhdr}
\usepackage{mathptmx} % Times New Roman
\usepackage{enumitem}\usepackage{amsfonts}
\usepackage{amsmath} % Recomendado para o comando \text{}
\usepackage{amssymb}
\usepackage{tikz}
\usetikzlibrary{positioning, shapes.geometric, calc}

% TikZ Styles for ER Diagrams
\tikzset{
    block/.style={rectangle, draw, thick, minimum width=2cm, minimum height=1cm, fill=white},
    diamond_node/.style={diamond, draw, thick, minimum width=2cm, minimum height=1cm, fill=white, aspect=2},
    attribute/.style={ellipse, draw, thin, minimum width=1.5cm, minimum height=0.6cm, fill=white, font=\scriptsize},
    line/.style={draw, thick},
}
% ------------------------------------------------ 
% Configuração de página
% ------------------------------------------------
\geometry{
    a4paper,
    left=2cm,
    right=2cm,
    top=2cm,
    bottom=2cm,
    headheight=15pt % Adicione isso
}% Caixa para resposta discursiva
\newcommand{\resposta}[1]{
\noindent
\framebox[\textwidth]{
\begin{minipage}{0.98\textwidth}
\vspace{#1}
\end{minipage}
}
}

\pagestyle{fancy}
\fancyhf{}
\renewcommand{\headrulewidth}{0pt}

\setlength{\parindent}{0pt}

% ------------------------------------------------
% Comandos auxiliares
% ------------------------------------------------
\newcolumntype{L}[1]{>{\raggedright\arraybackslash}p{#1}}
\newcolumntype{C}[1]{>{\centering\arraybackslash}p{#1}}

% ------------------------------------------------
% Documento
% ------------------------------------------------
\begin{document}

% ------------------------------------------------
% Cabeçalho - Corrigido para evitar aninhamento problemático
% ------------------------------------------------
\begin{tabularx}{\textwidth}{@{} L{4cm} X @{}}
    \raisebox{-0.5cm}{\includegraphics[width=4cm]{Universidade_de_Rio_Verde.jpg}} &
    \hfill
    \begin{tabular}{|>{\centering\arraybackslash}m{10cm}|>{\raggedright\arraybackslash}p{1.8cm}|}
        \hline
        \cellcolor{black}\color{white}\rule{0pt}{1.2cm}\textbf{Avaliação N1 (2ª CHAMADA)} 
        & 
        \scriptsize Nota \tabularnewline
        \hline
    \end{tabular} 
\end{tabularx}
\vspace{0.3cm}

% ------------------------------------------------  
% Tabela principal - Estrutura limpa
% ------------------------------------------------
\begin{tabularx}{\textwidth}{|X|L{2cm}|}
\hline
\textbf{Curso} & \textbf{Campus} \\
\textit{Engenharia de Software} & \textit{Rio Verde} \\
\hline
\multicolumn{2}{|l|}{\textbf{Disciplina}} \\  
\multicolumn{2}{|l|}{ \textit{Banco de dados}} \\
\hline
\textbf{Nome do(a) acadêmico(a)} & \textbf{Valor} \\
\textit{WAYSER HENRIQUE DE FARIA LOPES}& 7,0 \\
\hline
\textbf{Professor} & \textbf{Turma} \\
\textit{ Sergio Souza Novak} & \textit{B} \\
\hline
\textbf{Data da Avaliação} & \textbf{Código} \\
\textit{26/03/2026} & \textit{ESW417} \\
\hline
\end{tabularx}
 
% ------------------------------------------------
% Barra de Atenção
% ------------------------------------------------
\vspace{0.2cm}

\begin{tabularx}{\textwidth}{|X|}
\hline
\cellcolor{black}\color{white}
\centering \textbf{ATENÇÃO:} Somente serão passíveis de \textbf{REVISÃO} avaliações resolvidas à \textbf{TINTA}. \tabularnewline
\hline
\end{tabularx}

\vspace{0.3cm}

\textbf{ORIENTAÇÕES PARA A RESOLUÇÃO}  

\begin{itemize}
    \item A avaliação é individual, sem consulta;
    \item A interpretação de texto faz parte da avaliação;
\end{itemize}

\vspace{0.5cm}

\begin{center}
    \rule{\textwidth}{0.4pt}
    \textbf{QUESTÕES}
\end{center}

\vspace{0.5cm} 

\begin{enumerate}

 \item (ADAPTADO ENADE 2005) Todo funcionário deve estar alocado a uma única empresa prestadora de serviços. Assinale a opção que representa corretamente, no modelo entidade-relacionamento, a especificação apresentada acima.

\begin{enumerate}[label=\Alph*)]
    \item \begin{tikzpicture}[baseline=(rel.base), node distance=1cm, every node/.style={scale=0.8}]
        \node[block] (empresa) {empresa};
        \node[diamond_node, right=of empresa] (rel) {vínculo};
        \node[block, right=of rel] (func) {funcionário};
        \draw[line] (empresa) -- node[above] {(1,0)} node[below] {\scriptsize alocado em} (rel);
        \draw[line] (rel) -- node[above] {(n,0)} node[below] {\scriptsize possui} (func);
    \end{tikzpicture}
    
    \item \begin{tikzpicture}[baseline=(rel.base), node distance=1cm, every node/.style={scale=0.8}]
        \node[block] (empresa) {empresa};
        \node[diamond_node, right=of empresa] (rel) {vínculo};
        \node[block, right=of rel] (func) {funcionário};
        \draw[line] (empresa) -- node[above] {(0,1)} node[below] {\scriptsize alocado em} (rel);
        \draw[line] (rel) -- node[above] {(0,n)} node[below] {\scriptsize possui} (func);
    \end{tikzpicture}
    
    \item \begin{tikzpicture}[baseline=(rel.base), node distance=1cm, every node/.style={scale=0.8}]
        \node[block] (empresa) {empresa};
        \node[diamond_node, right=of empresa] (rel) {vínculo};
        \node[block, right=of rel] (func) {funcionário};
        \draw[line] (empresa) -- node[above] {(1,1)} node[below] {\scriptsize alocado em} (rel);
        \draw[line] (rel) -- node[above] {(0,n)} node[below] {\scriptsize possui} (func);
    \end{tikzpicture}
    \item \begin{tikzpicture}[baseline=(rel.base), node distance=1cm, every node/.style={scale=0.8}]
        \node[block] (empresa) {empresa};
        \node[diamond_node, right=of empresa] (rel) {vínculo};
        \node[block, right=of rel] (func) {funcionário};
        \draw[line] (empresa) -- node[above] {(1,n)} node[below] {\scriptsize alocado em} (rel);
        \draw[line] (rel) -- node[above] {(0,n)} node[below] {\scriptsize possui} (func);
    \end{tikzpicture}
     
    \item \begin{tikzpicture}[baseline=(rel.base), node distance=1cm, every node/.style={scale=0.8}]
        \node[block] (empresa) {empresa};
        \node[diamond_node, right=of empresa] (rel) {vínculo};
        \node[block, right=of rel] (func) {funcionário};
        \draw[line] (empresa) -- node[above] {(0,n)} node[below] {\scriptsize alocado em} (rel);
        \draw[line] (rel) -- node[above] {(0,n)} node[below] {\scriptsize possui} (func);
    \end{tikzpicture}
    
\end{enumerate}
    \item A correta alocação da Chave Estrangeira (FK) é fundamental para evitar a redundância de dados e o excesso de valores nulos (NULL). Com base nas definições de mapeamento lógico e físico (SQL), responda o que se pede. Uma locadora possui Filmes e Diretores. Um diretor pode dirigir vários filmes, mas cada filme é dirigido por apenas um único diretor. Assinale onde a FK deve ser inserida:
    \begin{enumerate}[label=\Alph*)]
        \item A FK deve ser inserida na tabela \textbf{Diretor}, pois é a entidade "forte".
        \item Deve-se criar uma terceira tabela para armazenar o relacionamento, por ser N:N.
        \item Ambos recebem FK um do outro para garantir a integridade.
        \item A FK deve ser inserida na tabela \textbf{Filme}, referenciando a Chave Primária de Diretor.
        \item Não é necessário o uso de FK neste cenário.
    \end{enumerate}

    \item \textbf{Implementação SQL:} Com base no diagrama lógico abaixo, que representa a tabela \textbf{ENGENHEIRO} em um mapeamento de especialização/generalização (Tabela Única), escreva o código SQL (\texttt{CREATE TABLE}) para criar esta tabela. Utilize um campo \textbf{id} como Chave Primária (SERIAL) e garanta que o campo \textbf{crea} seja único (UNIQUE). Apenas o campo departamento pode ser nulo.

    \begin{center}
    \begin{tikzpicture}[node distance=0cm, outer sep=0pt, every node/.style={draw, rectangle, minimum height=0.7cm, font=\small, anchor=north west}]
        \node (header) [fill=gray!20, minimum width=11cm] at (0,0) {\textbf{Tab.: ENGENHEIRO}};
        
        \node (col1) [minimum width=2cm] at (0,-0.7cm) {id (PK)};
        \node (col2) [minimum width=3cm] at (2cm,-0.7cm) {crea (UNIQUE)};
        \node (col3) [minimum width=3cm] at (5cm,-0.7cm) {nome};
        \node (col4) [minimum width=3cm] at (8cm,-0.7cm) {departamento};
        
        \node (r1c1) [minimum width=2cm] at (0,-1.4cm) {1};
        \node (r1c2) [minimum width=3cm] at (2cm,-1.4cm) {C-456};
        \node (r1c3) [minimum width=3cm] at (5cm,-1.4cm) {Eng. Carlos};
        \node (r1c4) [minimum width=3cm] at (8cm,-1.4cm) {Obras Civis};
        
        \node (r2c1) [minimum width=2cm] at (0,-2.1cm) {2};
        \node (r2c2) [minimum width=3cm] at (2cm,-2.1cm) {D-789};
        \node (r2c3) [minimum width=3cm] at (5cm,-2.1cm) {Eng. Diana};
        \node (r2c4) [minimum width=3cm, fill=yellow!10] at (8cm,-2.1cm) {NULL};
    \end{tikzpicture}
    \end{center}

    \resposta{5.5cm}

    \item \textbf{Relacionamento de Relacionamentos (Agregação):} O diagrama abaixo apresenta um \textbf{Relacionamento de Relacionamentos} (ou Agregação), onde o relacionamento \textbf{Venda} (entre Cliente e Vendedor) é associado a uma \textbf{Forma de Pagamento}. Realize a conversão desta estrutura para o Modelo Lógico, aplicando uma técnica que \textbf{elimine} a dependência do relacionamento de relacionamentos, transformando-o em uma estrutura de tabelas relacionais clássica (entidades e chaves estrangeiras).

    \begin{center}
    \begin{tikzpicture}[node distance=2.5cm, every node/.style={scale=0.85}]
        \node[block] (cli) {Cliente};
        \node[block, right=4cm of cli] (vend) {Vendedor};
        \node[diamond_node] (venda) at ($(cli)!0.5!(vend)$) {Venda};
        \draw[line] (cli) -- node[above] {(1,n)} (venda);
        \draw[line] (vend) -- node[above] {(1,n)} (venda);
        
        \draw[dashed, thick] ($(venda.north west)+(-0.6,0.6)$) rectangle ($(venda.south east)+(0.6,-0.6)$);
        \node[above=0.7cm of venda, font=\bfseries\small] {Agregação: Atendimento};

        \node[block, below=2.5cm of venda] (pag) {FormaPagto};
        \node[diamond_node, below=1cm of venda] (rel_pag) {utiliza};
        
        \draw[line] (venda) -- node[right] {(1,1)} (rel_pag);
        \draw[line] (rel_pag) -- node[right] {(1,n)} (pag);
        
        \node[attribute, left=0.5cm of cli] (c_pk) {\underline{id\_cli}};
        \draw (cli) -- (c_pk);
        \node[attribute, right=0.5cm of vend] (v_pk) {\underline{id\_vend}};
        \draw (vend) -- (v_pk);
        \node[attribute, right=0.5cm of pag] (p_pk) {\underline{id\_pg}};
        \draw (pag) -- (p_pk);
    \end{tikzpicture}
    \end{center}

    \resposta{8cm}
\item Considere o comando de criação da tabela abaixo. Assuma que o SGBD segue rigorosamente o padrão SQL e que nenhuma coluna possui valor padrão definido.

\begin{verbatim}
CREATE TABLE automoveis (
    id TEXT,
    motor DOUBLE PRECISION,
    portas INTEGER,
    ano INTEGER,
    preco NUMERIC,
    consumo DOUBLE PRECISION,
    torque DOUBLE PRECISION,
    potencia INTEGER,
    cilindrada INTEGER,
    velocidade_max INTEGER,
    aceleracao_0_100 DOUBLE PRECISION
);
\end{verbatim}

Analise as alternativas a seguir e assinale aquela que apresenta um comando \texttt{INSERT INTO} sintaticamente correto e semanticamente válido, considerando tipos de dados, correspondência de colunas e quantidade de valores.

\begin{enumerate}[label=\Alph*)]

    
    \item \texttt{INSERT INTO automoveis (id, motor, portas, ano, preco) VALUES ('C2', '2.0', 4, 2022, '75000');}
    \item \texttt{INSERT INTO automoveis VALUES ('C1', 2.0, 4, 2022, 75000, 12.5, 200, 150, 2000, 220, 8.5);}

    \item \texttt{INSERT INTO automoveis (motor, portas, ano, preco, consumo, torque, potencia, cilindrada, velocidade\_max, aceleracao\_0\_100) VALUES (2.0, 4, 2022, 75000, 12.5, 200, 150, 2000, 220, 8.5);}
    
    \item \texttt{INSERT INTO automoveis (id, motor, portas, ano, preco, consumo, torque, potencia, cilindrada, velocidade\_max, aceleracao\_0\_100) VALUES ('C4', 2.0, 4, 2022, 75000, 12.5, 200, 150, 2000, 220);}
    \item \texttt{INSERT INTO automoveis (id, motor, portas, ano, preco, consumo) VALUES ('C3', 2.0, 4, 2022, 75000, 12.5, 100);}


\end{enumerate}

\vspace{0.8cm}

    \item Um Autor pode escrever vários Livros, e um livro pode ser escrito por vários autores. Como este relacionamento é resolvido no modelo lógico?
    \begin{enumerate}[label=\Alph*)]
        \item Inserindo a FK de Autor na tabela \textbf{Livro}.
        \item Inserindo a FK de Livro na tabela \textbf{Autor}.
        \item Criando uma tabela associativa (ex: \textbf{Escritura}) que contém as FKs de ambas as tabelas originais.
        \item Mesclando as tabelas Autor e Livro em uma só.
        \item Transformando uma das entidades em um atributo multivalorado da outra.
    \end{enumerate}


\vspace{0.8cm}

 \item O diagrama a seguir apresenta um relacionamento ternário entre \textbf{Escritório}, \textbf{Analista} e \textbf{Projeto}. Com base nas regras de conversão do modelo conceitual para o modelo lógico, realize a \textbf{eliminação do ternário}, transformando-o em um conjunto de relacionamentos binários ou utilizando uma entidade associativa de forma a manter a integridade das regras de negócio.

\begin{center}
\begin{tikzpicture}[node distance=2.5cm, every node/.style={scale=0.9}]
    % Entidades
     \node[block] (escritorio) {Escritório};
    \node[block, below=of escritorio] (projeto) {Projeto};
    \node[block, left=of projeto, xshift=-1cm] (analista) {Analista};
    \node[block, right=of projeto, xshift=1cm] (equipo) {Equipamento};

    \node[attribute, left=0.5cm of escritorio] (esc_pk) {\underline{id\_esc}};
    \draw (escritorio) -- (esc_pk);
    
    \node[attribute, left=0.5cm of projeto] (pro_pk) {\underline{id\_proj}};
    \node[attribute, below left=0.5cm and 0.2cm of projeto] (pro_nome) {nome};
    \draw (projeto) -- (pro_pk);
    \draw (projeto) -- (pro_nome);
    
    \node[attribute, above=0.5cm of analista] (ana_pk) {\underline{id\_ana}};
    \node[attribute, above left=0.5cm and 0.5cm of analista] (ana_nome) {nome};
    \draw (analista) -- (ana_pk);
    \draw (analista) -- (ana_nome);
    
    \node[attribute, above=0.5cm of equipo] (eq_pk) {\underline{id\_eq}};
    \node[attribute, above right=0.5cm and 0.5cm of equipo] (eq_mod) {modelo};
    \draw (equipo) -- (eq_pk);
    \draw (equipo) -- (eq_mod);
    
    % Relacionamentos Binários
    \node[diamond_node, below=1cm of escritorio, yshift=0.5cm] (rel_local) {pertence};
    
    % Relacionamento Ternário
    \node[diamond_node, below=1.5cm of projeto] (rel_ternario) {alocação};
    \draw[line, <-] (rel_ternario) -- ++(3,-3) node[left] {Relacionamento Ternário}; 
    
    % Linhas e Cardinalidades
    \draw[line] (escritorio) -- node[right] {(1,1)} (rel_local);
    \draw[line] (rel_local) -- node[right] {(1,n)} (projeto);
    
    % Conexões do Ternário
    \draw[line] (projeto) -- node[left] {(1,n)} (rel_ternario);
    \draw[line] (analista) |- node[near start, left, yshift=0.5cm] {(1,n)} (rel_ternario);
    \draw[line] (equipo) |- node[near start, right, yshift=0.5cm] {(1,n)} (rel_ternario);
     
\end{tikzpicture}
\end{center}

\vspace{0.2cm}
\resposta{14cm}
\vspace{0.8cm}
\item Considere o diagrama lógico abaixo, representado segundo o modelo Entidade-Relacionamento Estendido (EER):

\begin{center}
\begin{tikzpicture}[node distance=1.2cm, every node/.style={draw, rectangle, minimum height=0.7cm, font=\small, fill=white}]
    % Superclasse
    \node (dispos) [minimum width=3cm] {\textbf{Dispositivo}};
    \node [below=0.1cm of dispos, draw=none, fill=none, font=\tiny] {(Superclasse)};
    
    % Símbolo de especialização 
    \node (circ) [circle, draw, below=1cm of dispos, inner sep=2pt] {S};
    \draw (dispos) -- (circ);
    
    % Subclasses
    \node (smart) [below left=1.2cm and 0.5cm of circ, minimum width=2.5cm] {\textbf{Smartphone}};
    \node (tablet) [below right=1.2cm and 0.5cm of circ, minimum width=2.5cm] {\textbf{Tablet}};
    
    \draw (circ) -| (smart);
    \draw (circ) -| (tablet);
    
    % Atributos
    \node (id) [draw, ellipse, above left=0.5cm of dispos] {\underline{id}};
    \draw (dispos.north west) -- (id);
    
    \node (os) [draw, ellipse, left=0.5cm of smart] {sis\_op};
    \draw (smart) -- (os);
    
    \node (tela) [draw, ellipse, right=0.5cm of tablet] {tam\_tela};
    \draw (tablet) -- (tela);
\end{tikzpicture}
\end{center}

Com base no diagrama apresentado, assinale a alternativa que descreve corretamente o conceito de modelagem utilizado \textbf{e} suas possíveis restrições semânticas no modelo EER.

\begin{enumerate}[label=\alph*)]

\item Trata-se de um caso de generalização total e disjunta, no qual toda entidade de \textbf{Dispositivo} deve pertencer exclusivamente a uma das subclasses.

\item Trata-se de uma especialização, podendo ser classificada como total ou parcial e como disjunta ou sobreposta, dependendo das restrições não explicitadas no diagrama.

\item Representa herança múltipla típica da orientação a objetos, na qual uma entidade pode herdar atributos simultaneamente de múltiplas superclasses.

\item Trata-se de um relacionamento identificador, no qual \textbf{Smartphone} e \textbf{Tablet} dependem existencialmente de \textbf{Dispositivo}.

\item Representa um caso de agregação, onde \textbf{Dispositivo} é composto por \textbf{Smartphone} e \textbf{Tablet} como partes integrantes.

\end{enumerate}

    \item Com base no diagrama ER abaixo, que representa um auto-relacionamento (recursivo) na entidade \textbf{Funcionario}, realize o mapeamento para o \textbf{Modelo Lógico Relacional}. Explique como a cardinalidade (1:N) influencia a alocação da Chave Estrangeira (FK) para representar a hierarquia de supervisão.

    \begin{center}
    \begin{tikzpicture}[node distance=2.5cm, every node/.style={scale=0.9}]
        \node[block] (func) {Funcionário};
        \node[diamond_node, right=of func] (ger) {gerencia};
        
        \draw[line] (func) -- node[above] {(0,1)} node[below] {\scriptsize supervisor} (ger);
        \draw[line] (ger) -- ++(0,-1.5) -| node[near start, below] {(0,n)} node[near end, left] {\scriptsize subordinados} (func);
        
        \node[attribute, left=1cm of func] (id) {\underline{id\_func}};
        \node[attribute, above left=0.5cm and 0.5cm of func] (nome) {nome};
        \draw (func) -- (id);
        \draw (func) -- (nome);
    \end{tikzpicture}
    \end{center}

    \resposta{10cm}


\end{enumerate}
\vspace{0.8cm}


\end{document} 
